We now extend the real-valued parameter-count statement of \Cref{cor:optimal-fact-storage-capacity} to a bounded-precision \emph{bit complexity} theorem.
The proof has three steps:
\textbf{(i)} positive margin implies robustness to output perturbations,
\textbf{(ii)} the bilinear MLP is Lipschitz in its parameters on the stored keys,
and \textbf{(iii)} sufficiently fine parameter quantization therefore preserves all margins.

\paragraph{Setup.}
Recall from \Cref{eq:mlp-construction} that our bilinear MLP has the form
\begin{equation}
g_{\boldsymbol{\theta}}(\mathbf{x})=\mathbf{B}\bigl((\mathbf{A}\mathbf{x})\odot(\mathbf{G}\mathbf{x})\bigr),
\qquad
\boldsymbol{\theta}=(\mathbf{A},\mathbf{G},\mathbf{B}),
\label{eq:bilinear_mlp_bits}
\end{equation}
with $\mathbf{A},\mathbf{G}\in\mathbb{R}^{m\times d}$ and $\mathbf{B}\in\mathbb{R}^{d\times m}$.
Let $P:=\dim(\boldsymbol{\theta})$ denote the total number of scalar parameters, so
$P\asymp md$ up to an absolute constant factor.

For a stored fact set $f:[F]\to[|V|]$, define the minimum margin
\[
\gamma_{\min}(\boldsymbol{\theta})
:=
\min_{i\in[F]}
\min_{j\neq f(i)}
\bigl\langle
g_{\boldsymbol{\theta}}(\mathbf{k}_{i}),\, \mathbf{v}_{f(i)}-\mathbf{v}_{j}
\bigr\rangle.
\]
We also write
\[
R_{K}:=\max_{i\in[F]}\|\mathbf{k}_{i}\|_{2},
\qquad
R_{V}:=\max_{a\in[|V|]}\|\mathbf{v}_{a}\|_{2}.
\]
Note that in the isotropic and unit-norm key/value regime of \Cref{sec:isotropic-case} and Appendix~\ref{app:theory:margin_bounds:rkrv}, one has $R_{K}\le 1$ and $R_{V}\le 1$.

\begin{lemma}[Margin robustness under output perturbations]
\label{lem:margin_robustness_bits}
Let $\boldsymbol{\theta}^\star$ be any parameter vector such that
$\gamma_{\min}(\boldsymbol{\theta}^\star)\ge \gamma_{0}>0$.
Assume that for some $\widetilde{\boldsymbol{\theta}}$,
\[
\max_{i\in[F]}
\bigl\|
g_{\widetilde{\boldsymbol{\theta}}}(\mathbf{k}_{i})-g_{\boldsymbol{\theta}^\star}(\mathbf{k}_{i})
\bigr\|_{2}
\le
\frac{\gamma_0}{4R_V}.
\]
Then $\widetilde{\boldsymbol{\theta}}$ stores the same fact set in the sense of
\Cref{def:fact-storage}, and in fact
\[
\gamma_{\min}(\widetilde{\boldsymbol{\theta}})\ge \frac{\gamma_0}{2}.
\]
\end{lemma}

\begin{proof}
Fix any stored key $\mathbf{k}_{i}$ and any competitor $j\neq f(i)$.
Then
\begin{align*}
\Bigl\langle
g_{\widetilde{\boldsymbol{\theta}}}(\mathbf{k}_{i}),\, \mathbf{v}_{f(i)}-\mathbf{v}_{j}
\Bigr\rangle
&=
\Bigl\langle
g_{\boldsymbol{\theta}^\star}(\mathbf{k}_{i}),\, \mathbf{v}_{f(i)}-\mathbf{v}_{j}
\Bigr\rangle
+
\Bigl\langle
g_{\widetilde{\boldsymbol{\theta}}}(\mathbf{k}_{i})-g_{\boldsymbol{\theta}^\star}(\mathbf{k}_{i}),\, \mathbf{v}_{f(i)}-\mathbf{v}_{j}
\Bigr\rangle \\
&\ge
\gamma_{0}
-
\bigl\|
g_{\widetilde{\boldsymbol{\theta}}}(\mathbf{k}_{i})-g_{\boldsymbol{\theta}^\star}(\mathbf{k}_{i})
\bigr\|_{2}
\,
\|\mathbf{v}_{f(i)}-\mathbf{v}_{j}\|_{2} \\
&\ge
\gamma_{0}
-
\frac{\gamma_0}{4R_V}
\bigl(\|\mathbf{v}_{f(i)}\|_{2}+\|\mathbf{v}_{j}\|_{2}\bigr) \\
&\ge
\gamma_{0}-\frac{\gamma_0}{4R_V}(2R_{V})
=
\frac{\gamma_0}{2}.
\end{align*}
Since this holds for every $i$ and every $j\neq f(i)$, we obtain $\gamma_{\min}(\widetilde{\boldsymbol{\theta}})\ge \gamma_{0}/2>0$, which implies that $\widetilde{\boldsymbol{\theta}}$ stores the same fact set.
\end{proof}

\begin{lemma}[Bilinear MLPs are Lipschitz in their parameters on stored keys]
\label{lem:bilinear_param_lipschitz_bits}
Fix any $\mathbf{x}\in\mathbb{R}^d$.
Let $\boldsymbol{\theta}=(\mathbf{A},\mathbf{G},\mathbf{B})$ and $\boldsymbol{\theta}'=(\mathbf{A}',\mathbf{G}',\mathbf{B}')$, and define
\[
M:=
\max\Bigl\{
\|\mathbf{A}\|_{\mathrm{op}},\|\mathbf{G}\|_{\mathrm{op}},\|\mathbf{B}\|_{\mathrm{op}},
\|\mathbf{A}'\|_{\mathrm{op}},\|\mathbf{G}'\|_{\mathrm{op}},\|\mathbf{B}'\|_{\mathrm{op}}
\Bigr\}.
\]
Then
\[
\bigl\|
g_{\boldsymbol{\theta}}(\mathbf{x})-g_{\boldsymbol{\theta}'}(\mathbf{x})
\bigr\|_{2}
\le
3M^2\|\mathbf{x}\|_{2}^2
\|\boldsymbol{\theta}-\boldsymbol{\theta}'\|_{2},
\]
where
\[
\|\boldsymbol{\theta}-\boldsymbol{\theta}'\|_{2}^2
:=
\|\mathbf{A}-\mathbf{A}'\|_{F}^2+\|\mathbf{G}-\mathbf{G}'\|_{F}^2+\|\mathbf{B}-\mathbf{B}'\|_{F}^2.
\]
Consequently, on any parameter region on which the operator norms of $\mathbf{A},\mathbf{G},\mathbf{B}$ are uniformly bounded by $M$, the map $\boldsymbol{\theta}\mapsto g_{\boldsymbol{\theta}}(\mathbf{k}_{i})$ is $3M^2R_{K}^2$-Lipschitz for every stored key $\mathbf{k}_{i}$ satisfying $\|\mathbf{k}_{i}\|_{2}\le R_{K}$.
\end{lemma}

\begin{proof}
Write
\begin{align*}
g_{\boldsymbol{\theta}}(\mathbf{x})-g_{\boldsymbol{\theta}'}(\mathbf{x})
&=
\mathbf{B}(\mathbf{A}\mathbf{x}\odot \mathbf{G}\mathbf{x})-\mathbf{B}'(\mathbf{A}'\mathbf{x}\odot \mathbf{G}'\mathbf{x}) \\
&=
(\mathbf{B}-\mathbf{B}')(\mathbf{A}\mathbf{x}\odot \mathbf{G}\mathbf{x})
+
\mathbf{B}'\Bigl((\mathbf{A}\mathbf{x}\odot \mathbf{G}\mathbf{x})-(\mathbf{A}'\mathbf{x}\odot \mathbf{G}'\mathbf{x})\Bigr) \\
&=
(\mathbf{B}-\mathbf{B}')(\mathbf{A}\mathbf{x}\odot \mathbf{G}\mathbf{x})
+
\mathbf{B}'\Bigl(((\mathbf{A}-\mathbf{A}')\mathbf{x})\odot \mathbf{G}\mathbf{x}\Bigr)
+
\mathbf{B}'\Bigl(\mathbf{A}'\mathbf{x}\odot ((\mathbf{G}-\mathbf{G}')\mathbf{x})\Bigr).
\end{align*}
Using
$\|\mathbf{u}\odot \mathbf{v}\|_{2} \le \|\mathbf{u}\|_{2}\|\mathbf{v}\|_\infty \le \|\mathbf{u}\|_{2}\|\mathbf{v}\|_{2}$
together with
$\|\mathbf{T}\mathbf{x}\|_{2}\le \|\mathbf{T}\|_{\mathrm{op}}\|\mathbf{x}\|_{2}$
gives
\begin{align*}
\|(\mathbf{B}-\mathbf{B}')(\mathbf{A}\mathbf{x}\odot \mathbf{G}\mathbf{x})\|_{2}
&\le
\|\mathbf{B}-\mathbf{B}'\|_{\mathrm{op}}\|\mathbf{A}\mathbf{x}\|_{2}\|\mathbf{G}\mathbf{x}\|_{2}
\le
M^2\|\mathbf{x}\|_{2}^2\|\mathbf{B}-\mathbf{B}'\|_{F},\\
\|\mathbf{B}'(((\mathbf{A}-\mathbf{A}')\mathbf{x})\odot \mathbf{G}\mathbf{x})\|_{2}
&\le
\|\mathbf{B}'\|_{\mathrm{op}}\|(\mathbf{A}-\mathbf{A}')\mathbf{x}\|_{2}\|\mathbf{G}\mathbf{x}\|_{2}
\le
M^2\|\mathbf{x}\|_{2}^2\|\mathbf{A}-\mathbf{A}'\|_{F},\\
\|\mathbf{B}'(\mathbf{A}'\mathbf{x}\odot ((\mathbf{G}-\mathbf{G}')\mathbf{x}))\|_{2}
&\le
\|\mathbf{B}'\|_{\mathrm{op}}\|\mathbf{A}'\mathbf{x}\|_{2}\|(\mathbf{G}-\mathbf{G}')\mathbf{x}\|_{2}
\le
M^2\|\mathbf{x}\|_{2}^2\|\mathbf{G}-\mathbf{G}'\|_{F}.
\end{align*}
Summing the three bounds and using
$a+b+c\le 3\sqrt{a^2+b^2+c^2}$ yields
\[
\|g_{\boldsymbol{\theta}}(\mathbf{x})-g_{\boldsymbol{\theta}'}(\mathbf{x})\|_{2}
\le
3M^2\|\mathbf{x}\|_{2}^2
\sqrt{
\|\mathbf{A}-\mathbf{A}'\|_{F}^2+\|\mathbf{G}-\mathbf{G}'\|_{F}^2+\|\mathbf{B}-\mathbf{B}'\|_{F}^2
}
=
3M^2\|\mathbf{x}\|_{2}^2\|\boldsymbol{\theta}-\boldsymbol{\theta}'\|_{2}.
\]
\end{proof}




\begin{theorem}[Bounded-bit implementation]
\label{thm:bounded_bit_implementation}
Fix a fact set and a real-valued parameter vector $\boldsymbol{\theta}^\star$ for a bilinear MLP of the form in \Cref{eq:bilinear_mlp_bits}.
Assume:
\begin{enumerate}
    \item[(i)] \textbf{Positive margin:}
    \[
    \gamma_{\min}(\boldsymbol{\theta}^\star)\ge \gamma_{0}>0.
    \]
    \item[(ii)] \textbf{Bounded dynamic range:}
    every coordinate of $\boldsymbol{\theta}^\star$ lies in $[-C_{\mathrm{rng}},C_{\mathrm{rng}}]$.
    \item[(iii)] \textbf{Parameter Lipschitzness on stored keys:}
    there exists $L_\theta>0$ such that for all
    $\boldsymbol{\theta},\boldsymbol{\theta}'\in[-C_{\mathrm{rng}},C_{\mathrm{rng}}]^P$,
    \begin{equation}
    \max_{i\in[F]}
    \|g_{\boldsymbol{\theta}}(\mathbf{k}_{i})-g_{\boldsymbol{\theta}'}(\mathbf{k}_{i})\|_{2}
    \le
    L_\theta\|\boldsymbol{\theta}-\boldsymbol{\theta}'\|_{2}.
    \label{eq:param_lipschitz_assumption_bits}
    \end{equation}
\end{enumerate}
Then there exists a floating-point parameter vector $\widetilde{\boldsymbol{\theta}}$
that stores the same fact set and satisfies
\[
\gamma_{\min}(\widetilde{\boldsymbol{\theta}})\ge \frac{\gamma_0}{2}.
\]
Moreover, one may choose $\widetilde{\boldsymbol{\theta}}$ by rounding each coordinate
of $\boldsymbol{\theta}^\star$ to a binary floating-point representation with $t$ mantissa bits, where
\[
t
=
O\!\left(
\log\!\left(
\frac{4C_{\mathrm{rng}}L_\theta R_{V}\sqrt{P}}{\gamma_0}
\right)
\right).
\]
In this case, each coordinate can be encoded using at most
\[
O\!\left(
\log\log C_{\mathrm{rng}}
+
\log\!\left(
\frac{4C_{\mathrm{rng}}L_\theta R_{V}\sqrt{P}}{\gamma_0}
\right)
\right)
\]
bits. Hence, the total number of bits needed to encode the quantized parameter vector
\(\widetilde{\boldsymbol{\theta}}\) is at most
\begin{equation}
\mathrm{Bits}(\widetilde{\boldsymbol{\theta}})
\le
P\cdot
O\!\left(
\log\log C_{\mathrm{rng}}
+
\log\!\left(
\frac{4C_{\mathrm{rng}}L_\theta R_{V}\sqrt{P}}{\gamma_0}
\right)
\right).
\label{eq:total_bits_bits}
\end{equation}
\end{theorem}

\begin{proof}
Let $\widetilde{\boldsymbol{\theta}}$ be obtained by rounding each coordinate of
$\boldsymbol{\theta}^\star$ to a binary floating-point representation with $t$ mantissa bits.
For standard floating-point rounding, each coordinate incurs relative error at most
$O(2^{-t})$, and hence absolute error at most
\[
O(2^{-t})\,C_{\mathrm{rng}}.
\]
Therefore
\[
\|\widetilde{\boldsymbol{\theta}}-\boldsymbol{\theta}^\star\|_{2}
\le
O(2^{-t})\,C_{\mathrm{rng}}\sqrt{P}.
\]
Choosing
\[
t
=
O\!\left(
\log\!\left(
\frac{4C_{\mathrm{rng}}L_\theta R_{V}\sqrt{P}}{\gamma_0}
\right)
\right)
\]
ensures that
\[
\|\widetilde{\boldsymbol{\theta}}-\boldsymbol{\theta}^\star\|_{2}
\le
\frac{\gamma_0}{4L_\theta R_V}.
\]
Applying the Lipschitz assumption (\Cref{eq:param_lipschitz_assumption_bits}),
\[
\max_{i\in[F]}
\|g_{\widetilde{\boldsymbol{\theta}}}(\mathbf{k}_{i})-g_{\boldsymbol{\theta}^\star}(\mathbf{k}_{i})\|_{2}
\le
L_\theta\|\widetilde{\boldsymbol{\theta}}-\boldsymbol{\theta}^\star\|_{2}
\le
\frac{\gamma_0}{4R_V}.
\]
Lemma~\ref{lem:margin_robustness_bits} therefore implies that
$\widetilde{\boldsymbol{\theta}}$ stores the same fact set and that
$\gamma_{\min}(\widetilde{\boldsymbol{\theta}})\ge \gamma_{0}/2$.

It remains to count bits.
Each nonzero floating-point coordinate can be written in the form
\[
\pm m\,2^e,
\]
where the mantissa $m$ is represented to $t$ bits of precision.
Since every coordinate lies in $[-C_{\mathrm{rng}},C_{\mathrm{rng}}]$, the exponent satisfies
$|e|=O(\log C_{\mathrm{rng}})$, and hence the exponent can be encoded using
$O(\log\log C_{\mathrm{rng}})$ bits.
Thus each coordinate can be encoded using at most
\[
O\!\left(
\log\log C_{\mathrm{rng}}
+
t
\right)
=
O\!\left(
\log\log C_{\mathrm{rng}}
+
\log\!\left(
\frac{4C_{\mathrm{rng}}L_\theta R_{V}\sqrt{P}}{\gamma_0}
\right)
\right)
\]
bits.
Multiplying by the \(P\) coordinates yields~\Cref{eq:total_bits_bits}.
\end{proof}

Next, we specialize the previous bit complexity theorem to our sketched bilinear MLP construction (\Cref{alg:sketched_k2_hebbian_whitening}) by bounding its dynamic range and Lipschitz constant. In the isotropic unit-norm regime of \Cref{sec:isotropic-case} and Appendix~\ref{app:theory:margin_bounds:rkrv}, they reduce to $C_{\mathrm{rng}},L_\theta\le d^{O(1)}$, yielding an explicit bit-complexity bound.

\begin{proposition}[Dynamic range and stored-key Lipschitzness for the unwhitened sketched-$K_{2}$ construction]
\label{prop:unwhitened_k2_norm_bounds}
Let \(\mathbf{a}_1,\dots,\mathbf{a}_m,\mathbf{b}_1,\dots,\mathbf{b}_m \in \mathbb{R}^d\) be sampled i.i.d. from \(N(0,\mathbf{I}_d)\), and consider the sketched-$K_{2}$ feature map
\[
g(\mathbf{x})=\frac{1}{\sqrt m}\bigl((\mathbf{a}_r^\top \mathbf{x})(\mathbf{b}_r^\top \mathbf{x})\bigr)_{r=1}^m.
\]
Let \(\bar{\mathbf{A}},\bar{\mathbf{G}}\in\mathbb{R}^{m\times d}\) be the matrices whose rows are \(\mathbf{a}_r^\top\) and \(\mathbf{b}_r^\top\), respectively, and realize \(g\) in bilinear MLP form by setting
\[
\mathbf{A}=m^{-1/4}\bar{\mathbf{A}},
\qquad
\mathbf{G}=m^{-1/4}\bar{\mathbf{G}}.
\]
Let the raw Hebbian readout be
\[
\mathbf{B}_0=\frac{1}{F}\mathbf{C}_f^\top \boldsymbol{\Phi},
\]
where the rows of \(\boldsymbol{\Phi}\in\mathbb{R}^{F\times m}\) are \(g(\mathbf{k}_i)^\top\)
and the rows of \(\mathbf{C}_f\in\mathbb{R}^{F\times d}\) are the value embeddings \(\mathbf{v}_{f(i)}^\top\).

Assume
\[
m,\;F,\;\delta^{-1}\le d^{O(1)}.
\]
Then, with probability at least \(1-\delta\) over the draw of
\(\{\mathbf{a}_r,\mathbf{b}_r\}_{r=1}^m\), the resulting bilinear MLP
\[
\mathbf{x}\mapsto \mathbf{B}_{0}\bigl((\mathbf{A}\mathbf{x})\odot(\mathbf{G}\mathbf{x})\bigr)
\]
has parameter dynamic range and stored-key parameter Lipschitz constant bounded by
\[
C_{\mathrm{rng}},\;L_\theta \le d^{O(1)}\,\poly(R_{K},R_{V}).
\]
where one may take
\[
C_{\mathrm{rng}} \le d^{O(1)}\max\{1,R_VR_{K}^2\},
\qquad
L_\theta \le d^{O(1)}\,R_{K}^2\max\{1,R_{V}^2R_{K}^4\}.
\]
\end{proposition}

\begin{proof}
Standard Gaussian random matrix bounds~\citep{vershynin2018high} imply that with probability at least $1-\delta$ over the draw of \(\{\mathbf{a}_r,\mathbf{b}_r\}_{r=1}^m\),
\[
\|\mathbf{A}\|_{\mathrm{op}},\|\mathbf{G}\|_{\mathrm{op}}
\lesssim
m^{-1/4}\Bigl(\sqrt m+\sqrt d+\sqrt{\log(1/\delta)}\Bigr).
\]
Under the assumption $m,\delta^{-1}\le d^{O(1)}$, it follows that
\[
\|\mathbf{A}\|_{\mathrm{op}},\|\mathbf{G}\|_{\mathrm{op}} \le d^{O(1)}.
\]

Next, since each value embedding has norm at most $R_{V}$,
\[
\|\mathbf{C}_{f}\|_{\mathrm{op}}\le \|\mathbf{C}_{f}\|_{F} \le \sqrt{F}\,R_{V}.
\]
Also,
\[
\|\boldsymbol{\Phi}\|_{\mathrm{op}}
\le
\|\boldsymbol{\Phi}\|_{F}
\le
\sqrt{F}\,\max_{i\in[F]}\|g(\mathbf{k}_{i})\|_{2}.
\]
For every stored key $\mathbf{k}_{i}$,
\[
\|g(\mathbf{k}_{i})\|_{2}
=
\bigl\|(\mathbf{A}\mathbf{k}_{i})\odot(\mathbf{G}\mathbf{k}_{i})\bigr\|_{2}
\le
\|\mathbf{A}\mathbf{k}_{i}\|_{2}\,\|\mathbf{G}\mathbf{k}_{i}\|_{2}
\le
\|\mathbf{A}\|_{\mathrm{op}}\|\mathbf{G}\|_{\mathrm{op}}R_{K}^2.
\]
Therefore
\[
\|\boldsymbol{\Phi}\|_{\mathrm{op}}
\le
\sqrt{F}\,\|\mathbf{A}\|_{\mathrm{op}}\|\mathbf{G}\|_{\mathrm{op}}R_{K}^2,
\]
and hence
\[
\|\mathbf{B}_{0}\|_{\mathrm{op}}
\le
\frac{1}{F}\|\mathbf{C}_{f}\|_{\mathrm{op}}\|\boldsymbol{\Phi}\|_{\mathrm{op}}
\le
R_{V}\|\mathbf{A}\|_{\mathrm{op}}\|\mathbf{G}\|_{\mathrm{op}}R_{K}^2
\le
d^{O(1)}R_VR_{K}^2
\]
with probability at least \(1-\delta\).

Now define
\[
M:=\max\{\|\mathbf{A}\|_{\mathrm{op}},\|\mathbf{G}\|_{\mathrm{op}},\|\mathbf{B}_{0}\|_{\mathrm{op}}\}.
\]
From the previous bounds,
\[
M \le d^{O(1)}\max\{1,R_VR_{K}^2\}
\]
with probability at least \(1-\delta\).

To pass from operator norms to coordinate bounds, note that for any matrix $\mathbf{T}$,
\[
|T_{ab}| = |\mathbf{e}_{a}^\top \mathbf{T} \mathbf{e}_{b}| \le \|\mathbf{T}\|_{\mathrm{op}}.
\]
Hence every scalar entry of $\mathbf{A},\mathbf{G},\mathbf{B}_{0}$ is bounded by $M$, and so one may take
\[
C_{\mathrm{rng}}\le M \le d^{O(1)}\max\{1,R_VR_{K}^2\}.
\]

Finally, \Cref{lem:bilinear_param_lipschitz_bits} yields
\[
L_\theta \le 3M^2R_{K}^2
\le
d^{O(1)}\,R_{K}^2\max\{1,R_{V}^2R_{K}^4\},
\]
which proves the claim.
\end{proof}



\begin{corollary}[Bit complexity in the isotropic regime for the unwhitened construction]
\label{cor:bit_complexity_isotropic}
Assume the isotropic key/value regime of
\Cref{cor:optimal-fact-storage-capacity}, and consider the sketched-$K_{2}$ construction of~\Cref{alg:sketched_k2_hebbian_whitening}.
Choose the width $m$ so that the corresponding real-valued construction
satisfies
\[
P\asymp md \asymp F\log(F/\delta)
\]
and has constant slack margin
\[
\gamma_{\min}(\boldsymbol{\theta}^\star)\ge c_{0}>0
\]
with high probability.
Assume further that $F,\delta^{-1}\le d^C$ for an absolute constant $C$.
Then with high probability the same fact set can be stored using
\begin{equation}
\mathrm{Bits}(\widetilde{\boldsymbol{\theta}})
=
O\!\Bigl(
F\log(F/\delta)\log d
\Bigr).
\label{eq:bit_complexity_isotropic_bits}
\end{equation}
\end{corollary}

\begin{proof}
By the isotropic real-valued capacity result, the unwhitened construction stores
$F$ facts using
\[
P\asymp md\asymp F\log(F/\delta)
\]
parameters, up to an absolute constant factor.
In the isotropic unit-norm regime, $R_{K},R_{V}\le 1$.
Since $F \leq d^{O(1)}$, \Cref{prop:unwhitened_k2_norm_bounds} yields
\[
C_{\mathrm{rng}},L_\theta \le d^{O(1)}
\]
with high probability.
Applying \Cref{thm:bounded_bit_implementation},
\[
\mathrm{Bits}(\widetilde{\boldsymbol{\theta}})
\le
P\,
O\!\left(
\log\log C_{\mathrm{rng}}
+
\log\!\left(
\frac{4C_{\mathrm{rng}}L_\theta R_{V}\sqrt{P}}{c_0}
\right)
\right).
\]
Since $C_{\mathrm{rng}},L_\theta,P\le d^{O(1)}$ under the standing assumption
$F,\delta^{-1}\le d^C$, the quantity inside the outer \(O(\cdot)\) is \(O(\log d)\).
Combining this with $P\asymp F\log(F/\delta)$ yields
\Cref{eq:bit_complexity_isotropic_bits}.
\end{proof}
